\section*{The Propers: Themes and Movement}

The solemnity moves from remembered dependence to sacramental participation and from participation to unity, adoration, and hope. Its printed alternatives give distinct textual paths through the same celebration, each on the terms the governing books supply.

\begin{center}
\begin{tikzpicture}[x=3.35cm]
\node[guide node,text width=2.55cm] (memory) at (0,0) {\textbf{Dependence remembered}\\Wilderness, word, manna};
\node[guide node,text width=2.55cm] (gift) at (1,0) {\textbf{Gift fulfilled}\\Christ, true food and drink};
\node[guide emphasis,text width=2.55cm] (body) at (2,0) {\textbf{Communion forms}\\One bread, one body};
\node[guide node,text width=2.55cm] (hope) at (3,0) {\textbf{Life extends}\\Unity, adoration, glory};
\foreach \x/\y in {memory/gift,gift/body,body/hope}
  \draw[guide arrow] (\x) -- (\y);
\end{tikzpicture}
\end{center}

\subsection*{Remembering Dependence\enspace\elementref{Entrance, Collect, First Reading, Responsorial Psalm}}

The Entrance adapts Psalm 81's covenantal promise; the Collect names Eucharist as memorial of the Passion; Deuteronomy warns against forgetting the God who fed and led; Psalm 147 praises the God whose word and wheat restore Jerusalem. John Paul II directly joins Deuteronomy's imperative of memory to Eucharistic anamnesis. The joining is John Paul II's; the later sacramental vocabulary is not that of Deuteronomy's historical author.

\begin{sourcecard}{John Paul II, Corpus Christi homily, 2002}
John Paul II places Deuteronomy's “remember” and “do not forget” beside the Eucharistic memorial of Christ's saving Passion. The homily supports this formula's movement from wilderness dependence to sacramental remembrance without erasing either text's historical setting.
\end{sourcecard}

\subsection*{Participation Makes One Body\enspace\elementref{Second Reading, Prayer over Offerings}}

Paul's cup and bread establish real participation, one ecclesial body, and a boundary against idolatry. The Prayer over the Offerings asks for unity and peace signified by the gifts. Augustine's direct conjunction of John 6 and 1 Corinthians 10 gives the inherited phrase “sign of unity, bond of charity” its demanding force: a good sacrament may be received badly, so sacramental participation calls for reconciliation rather than sentimental claims of harmony.

\subsection*{The Living Bread Gives Abiding Life\enspace\elementref{Sequence, Acclamation, Gospel, Prefaces}}

The sequence, when used, gives poetic catechesis before the Alleluia. John then proclaims flesh and blood, mutual abiding, eternal life, and last-day resurrection. Preface I places the gift inside Christ's priesthood and the one saving sacrifice; Preface II gathers memorial, love, unity, banquet, grace, and pledge. The two Prefaces are alternative liturgical lenses, not competing doctrines, and the Gospel's meaning does not depend on the optional sequence.

\begin{studybox}{Real presence in the one Eucharist}
The Church confesses sacramental, true, real, and substantial presence under the species---neither a bare mental symbol nor crude materialism. The one Eucharist is sacrifice, memorial, presence, banquet, ecclesial communion, and pledge. Christ is not divided, and saving fruit is not mechanically guaranteed apart from worthy reception and conversion.
\end{studybox}

\subsection*{Reception Opens toward Glory\enspace\elementref{Communion, Prayer after Communion}}

Communion centers mutual abiding; the final prayer stretches present reception toward eternal participation in divine life. The sacramental gift thus joins present communion, ecclesial unity, bodily resurrection, and the hope of glory.
